%%%%%%%%%%%%%%%%%
\begin{figure}[!t]
    \centering
    \includegraphics[width=0.95\columnwidth]{figures/results/track_loop.pdf}
    \caption{Camera-pose tracking for the loop: position error (top) stays sub-metre while the camera orientation error (bottom, log scale) remains below $0.01^\circ$ throughout, despite the airframe inverting.}
    \label{fig:track_loop}
\end{figure}
%%%%%%%%%%%%%%%%%

\section{Validation Results}
\label{sec:results}

We validate the proposed framework on a representative $\SI{5}{\kilo\gram}$ multicopter equipped with a 3-axis gimbal, commanded to reproduce the camera pose of a Red Bull Edge~540~\cite{redbull} aerobatic aircraft. The Edge~540 is chosen for its aggressive maneuvering envelope, providing a demanding reference. Vehicle parameters are listed in Table~\ref{tab:vehicle_params}. The reference trajectories are generated by the full $6$-DOF fixed-wing model~\eqref{eq:plane_dynamics}, which supplies the camera position and its derivatives up to snap, and the camera attitude, as the inputs $\mathbf{v}_\text{MC}$ to the dynamic-feedback-linearizing law~\eqref{eq:feedback_control_extended}.

\begin{table}[!b]
\centering
\caption{Vehicle Parameters}
\label{tab:vehicle_params}
\begin{tabular}{|l|l|}
\hline
\multicolumn{2}{|c|}{\textbf{Fixed-Wing (Red Bull Edge 540)}} \\
\hline
Mass & 530 kg \\
Moments of inertia & (2419, 2133, 4552) kg·m² \\
Wingspan & 7.40 m \\
Wing Area & 9.10 m² \\
Max Speed & 118 m/s \\
\hline
\multicolumn{2}{|c|}{\textbf{Multicopter (3-axis gimbal)}} \\
\hline
Mass & 5.0 kg \\
Moments of inertia & (0.15, 0.15, 0.25) kg·m² \\
Arm length & 0.35 m \\
\hline
\end{tabular}%
\end{table}

%%%%%%%%%%%%%%%%%%%%%%%%%%%%%%%%%%%%%%%%%%%%%%%%%%%%%%%%%%%%%%%%%%%%%%%%%%%%%%%%
\subsection{Controller Realization}
\label{ssec:realization}
The exact linearization of Section~\ref{sec:method} reduces each output to an integrator chain $y_i^{(r_i)} = v_i$. We close these chains with three theory-grounded elements and \emph{no} ad-hoc clamps:
\begin{itemize}[leftmargin=*]
\item \emph{One-parameter critically-damped gains.} All poles of every error chain are placed at a single bandwidth $-\omega$, so the feedback gains are the binomial coefficients of $(s+\omega)^{r_i}$. There is no hand-tuned gain vector; we use $\omega = \SI{30}{\radian\per\second}$.
\item \emph{Singularity-robust inverse.} The decoupling map carries a $1/q_0$ term that diverges at the body-attitude representation singularity ($q_0 \to 0$, i.e.\ the airframe inverting past $90^\circ$). We floor $|q_0|$ when evaluating the inverse, keeping the linearizing law bounded through inverted flight while remaining exact away from the singular set.
\item \emph{Reference governor.} Since the thrust and rate demands of a maneuver are set by the reference and are mass-invariant under feedback linearization, we admit a pace factor $\lambda \in (0,1]$ that time-scales the reference so the inverse inputs stay within the actuator and gimbal limits, i.e.\ within the feasible set~$\mathcal{F}$.
\end{itemize}

%%%%%%%%%%%%%%%%%%%%%%%%%%%%%%%%%%%%%%%%%%%%%%%%%%%%%%%%%%%%%%%%%%%%%%%%%%%%%%%%
\subsection{Tracking Performance}
\label{ssec:tracking}
We evaluate six aerobatic maneuvers: a vertical loop, an axial roll, a soft roll, a barrel roll, a Split-S, and a Cuban-8. Each reference is paced by the governor to a feasible thrust-to-weight ($\text{T/W} \le 3$). Table~\ref{tab:results} reports, per maneuver, the governor pace $\lambda$, the peak thrust-to-weight, the mean-absolute camera-position error, and the maximum camera-orientation (geodesic) error.

\begin{table}[!t]
\centering
\caption{Camera-Pose Tracking (governed, $\text{T/W} \le 3$)}
\label{tab:results}
\resizebox{\columnwidth}{!}{%
\begin{tabular}{|l|c|c|c|c|}
\hline
\textbf{Maneuver} & $\lambda$ & peak T/W & MAE pos. (m) & max orient. (deg) \\
\hline
Loop        & 0.87 & 3.0 & 0.14 & 0.003 \\
Roll        & 0.60 & 3.0 & 0.07 & 0.006 \\
Soft roll   & 0.60 & 3.0 & 0.11 & 0.021 \\
Barrel roll & 0.60 & 3.0 & 0.10 & 0.028 \\
Split-S     & 0.29 & 3.0 & 0.06 & 0.003 \\
Cuban-8     & 0.35 & 3.0 & 0.07 & 0.008 \\
\hline
\end{tabular}}
\end{table}

Across all maneuvers the camera-orientation error stays below $0.05^\circ$ and the camera-position error stays sub-metre, confirming high-fidelity reproduction of the fixed-wing camera pose. Fig.~\ref{fig:track_loop} shows the loop in detail: both the position error and the (log-scale) orientation error remain small throughout, even as the airframe passes through inverted flight.

Fig.~\ref{fig:overlays} overlays, in three dimensions, the fixed-wing camera path (red) and the multicopter camera path (blue) for four maneuvers, with the (transparent) aircraft drawn at intervals along with the multicopter camera frame (green boresight). The two paths are visually indistinguishable, and the camera boresight stays aligned with the aircraft, confirming that position and orientation are both reproduced. The figure also makes the decoupling explicit: the airframe attitude (transparent aircraft pose) varies strongly along the trajectory while the camera frame remains locked on the reference.

\begin{figure}[!t]
    \centering
    \subfloat[Loop]{\includegraphics[width=0.48\columnwidth]{figures/results/overlay_loop.pdf}}\hfill
    \subfloat[Roll]{\includegraphics[width=0.48\columnwidth]{figures/results/overlay_roll.pdf}}\\
    \subfloat[Split-S]{\includegraphics[width=0.48\columnwidth]{figures/results/overlay_splits.pdf}}\hfill
    \subfloat[Cuban-8]{\includegraphics[width=0.48\columnwidth]{figures/results/overlay_cuban8.pdf}}
    \caption{3D camera-pose overlays. Fixed-wing camera path (red) and multicopter camera path (blue, dashed) coincide; the transparent aircraft and the multicopter camera frame (green boresight) are shown at intervals. The airframe attitude varies while the camera stays locked on the reference.}
    \label{fig:overlays}
\end{figure}

The decoupling is the key behaviour. To match the fixed-wing \emph{position}, the multicopter must reproduce its specific-force vector; the thrust direction tracks the aircraft's specific force to within ${\sim}10^\circ$, so the airframe tilts and even inverts (body tilt reaching $170$--$180^\circ$) to vector thrust along the curved, high-load trajectory. The 3-axis gimbal simultaneously absorbs the residual rotation $\mathbf{q}_\text{G} = \bar{\mathbf{q}}_\text{M} \otimes \mathbf{q}_\text{A}$, so the \emph{camera} orientation stays locked on the fixed-wing's. The airframe thus reproduces the load factor experienced by the fixed-wing while the camera viewpoint remains faithful, which is the goal of physical-to-physical emulation.

%%%%%%%%%%%%%%%%%%%%%%%%%%%%%%%%%%%%%%%%%%%%%%%%%%%%%%%%%%%%%%%%%%%%%%%%%%%%%%%%
\subsection{Control Inputs and Feasibility}
\label{ssec:feasibility}
Fig.~\ref{fig:inputs} shows the realized control effort for the loop and roll against the platform limits (thrust-to-weight $\le 3$, body rate $\le \SI{600}{\deg\per\second}$, gimbal rate $\le \SI{400}{\deg\per\second}$). After a brief initial transient, all three inputs remain within their limits for the full maneuver: thrust-to-weight oscillates between roughly $1$ and $3$ as the aircraft loads and unloads, while the body and gimbal rates settle well below their bounds.

\begin{figure}[!t]
    \centering
    \subfloat[Loop]{\includegraphics[width=0.48\columnwidth]{figures/results/inputs_loop.pdf}}\hfill
    \subfloat[Roll]{\includegraphics[width=0.48\columnwidth]{figures/results/inputs_roll.pdf}}
    \caption{Control inputs versus feasibility limits (shaded = infeasible). Thrust-to-weight, body rate, and gimbal rate all remain within bounds after the initial transient.}
    \label{fig:inputs}
\end{figure}

Because the exact inverse reproduces the reference acceleration regardless of vehicle mass, the thrust-to-weight demand is a property of the \emph{trajectory}, not the platform. At full speed the Edge~540 maneuvers demand $4$--$39\times$ thrust-to-weight (Fig.~\ref{fig:feasibility}), far beyond any realizable multicopter. The reference governor pulls each maneuver into $\mathcal{F}$ (Lemma~\ref{lem:feasible}) by pacing, recovering $\text{T/W} \le 3$ while preserving exact tracking (Table~\ref{tab:results}).

\begin{figure}[!t]
    \centering
    \includegraphics[width=0.92\columnwidth]{figures/results/feasibility.pdf}
    \caption{Peak thrust-to-weight per maneuver at full speed versus governed; the governor confines the demand below the platform limit ($\text{T/W}\le 3$).}
    \label{fig:feasibility}
\end{figure}

Feasibility is bounded \emph{jointly} by thrust authority and the gimbal kinematic range. Slowing a maneuver reduces the thrust demand but also reduces the airframe tilt, shifting orientation onto the gimbal and toward its $\theta_\text{G} = \pm\frac{\pi}{2}$ singularity. The loop and the axial roll satisfy both limits comfortably (gimbal margin $\cos\theta_\text{G} \ge 0.38$); the remaining maneuvers track exactly and respect thrust but operate nearer the gimbal limit. This interplay, in which the body flip extends the orientation range beyond the gimbal at the cost of thrust authority, characterizes the boundary of $\mathcal{F}$ for a given platform.

%%%%%%%%%%%%%%%%%%%%%%%%%%%%%%%%%%%%%%%%%%%%%%%%%%%%%%%%%%%%%%%%%%%%%%%%%%%%%%%%
\subsection{Discussion}
The results validate the framework of Sections~\ref{sec:proof} and~\ref{sec:method}: the multicopter--gimbal system achieves camera-pose mimicry $\mathbf{y}_\text{MC}(t) \equiv \mathbf{y}_\text{AC}(t)$ with sub-degree orientation error and sub-metre position error across a suite of aerobatic maneuvers, using an exact, singularity-robust feedback-linearizing controller with a single stabilization parameter. The reference governor renders any reference feasible by confining it to $\mathcal{F}$ (Lemma~\ref{lem:feasible}), replacing the non-constructive Assumption~\ref{ass:authority} with a checkable condition derived directly from the system inverse.
